\documentclass[aspectratio=169,10pt]{beamer}

\usepackage{fontspec}
\usepackage{polyglossia}
\setdefaultlanguage{vietnamese}
\setotherlanguage{english}
\setmainfont{DejaVu Sans}
\setsansfont{DejaVu Sans}
\setmonofont{DejaVu Sans Mono}

\usepackage{tikz}
\usetikzlibrary{arrows.meta,positioning,calc}
\usepackage{xcolor}
\usepackage{listings}
\usepackage{amsmath}
\usepackage{booktabs}

\usetheme{Madrid}
\usecolortheme{seahorse}
\setbeamertemplate{navigation symbols}{}
\newcommand{\TotalFramesFixed}{42}
\setbeamertemplate{footline}{%
  \leavevmode%
  \hbox{\begin{beamercolorbox}[wd=\paperwidth,ht=2.6ex,dp=1.1ex,right]{page number in head/foot}%
    \usebeamerfont{page number in head/foot}\insertframenumber/\TotalFramesFixed\hspace*{1.5ex}
  \end{beamercolorbox}}%
}

\definecolor{pyblue}{HTML}{1F77B4}
\definecolor{pygreen}{HTML}{2CA02C}
\definecolor{pyorange}{HTML}{FF7F0E}
\definecolor{pypurple}{HTML}{9467BD}
\definecolor{pygray}{HTML}{F2F4F8}
\definecolor{dark}{HTML}{1F2937}

\lstdefinestyle{pythonstyle}{
  language=Python,
  basicstyle=\ttfamily\small,
  keywordstyle=\color{pyblue}\bfseries,
  commentstyle=\color{pygreen},
  stringstyle=\color{pyorange},
  showstringspaces=false,
  columns=fullflexible,
  keepspaces=true,
  frame=single,
  rulecolor=\color{gray!40},
  backgroundcolor=\color{pygray},
  breaklines=true
}
\lstset{style=pythonstyle}

\tikzset{
  box/.style={draw=dark, rounded corners=2mm, very thick, align=center, fill=pygray, minimum height=10mm, inner sep=2.5mm},
  bluebox/.style={box, fill=blue!12},
  greenbox/.style={box, fill=green!12},
  orangebox/.style={box, fill=orange!16},
  purplebox/.style={box, fill=purple!12},
  redbox/.style={box, fill=red!12},
  arrow/.style={-{Latex[length=2.5mm]}, very thick, draw=dark}
}

\title[Nhập/Xuất cơ bản]{Nhập/Xuất cơ bản trong Python}
\subtitle{Sử dụng \texttt{print()}, \texttt{input()}, chuyển kiểu và f-string}
\author{Nhập môn lập trình Python}
\date{}

\begin{document}

\begin{frame}
  \titlepage
\end{frame}

\begin{frame}{Cấu trúc bài học}
\begin{enumerate}
  \item Xuất dữ liệu với \texttt{print()}
  \item Nhập dữ liệu từ bàn phím với \texttt{input()}
  \item Chuyển đổi kiểu dữ liệu đầu vào
  \item Mô hình Input $\rightarrow$ Process $\rightarrow$ Output
  \item Định dạng output với f-string
  \item Các lỗi nhập/xuất thường gặp
  \item Bài tập thực hành tổng hợp
\end{enumerate}
\end{frame}

\begin{frame}{Mục tiêu học tập}
Sau bài học này, sinh viên có thể:
\begin{itemize}
  \item sử dụng \texttt{print()} để xuất dữ liệu;
  \item in nhiều giá trị với \texttt{sep} và \texttt{end};
  \item sử dụng \texttt{input()} để nhận dữ liệu từ bàn phím;
  \item giải thích vì sao \texttt{input()} trả về \texttt{str};
  \item chuyển input sang \texttt{int} hoặc \texttt{float};
  \item kết hợp nhập, xử lý và xuất trong một chương trình;
  \item định dạng output bằng f-string cơ bản.
\end{itemize}
\begin{alertblock}{Phạm vi bài học}
Tất cả chương trình được viết trực tiếp bằng các câu lệnh cơ bản. Chưa sử dụng hàm tự định nghĩa (\texttt{def}).
\end{alertblock}
\end{frame}

\begin{frame}{Mẫu chương trình cơ bản}
\begin{center}
\begin{tikzpicture}[scale=0.95, every node/.style={transform shape}]
  \node[bluebox, minimum width=3.2cm] (input) at (0,0) {Input\\nhận dữ liệu};
  \node[orangebox, minimum width=3.2cm] (process) at (4.5,0) {Process\\tính toán / biến đổi};
  \node[greenbox, minimum width=3.2cm] (output) at (9.0,0) {Output\\hiển thị kết quả};
  \draw[arrow] (input.east) -- (process.west);
  \draw[arrow] (process.east) -- (output.west);
\end{tikzpicture}
\end{center}
\vspace{0.3cm}
\begin{block}{Ví dụ: diện tích hình chữ nhật}
Nhập chiều dài và chiều rộng $\rightarrow$ tính diện tích $\rightarrow$ in kết quả.
\end{block}
\end{frame}

\section{Xuất dữ liệu với print()}

\begin{frame}[fragile]{Xuất dữ liệu với \texttt{print()}}
\begin{block}{Cú pháp cơ bản}
\begin{lstlisting}
print(value)
\end{lstlisting}
\end{block}
\begin{exampleblock}{Ví dụ}
\begin{lstlisting}
print("Hello, Python!")
\end{lstlisting}
\end{exampleblock}
\begin{block}{Output}
\begin{verbatim}
Hello, Python!
\end{verbatim}
\end{block}
\end{frame}

\begin{frame}[fragile]{In số và biến}
\begin{columns}[T,onlytextwidth]
\column{0.48\textwidth}
\begin{block}{Giá trị số}
\begin{lstlisting}
print(100)
print(3.14)
\end{lstlisting}
\end{block}
\column{0.48\textwidth}
\begin{block}{Biến}
\begin{lstlisting}
name = "An"
age = 18
print(name)
print(age)
\end{lstlisting}
\end{block}
\end{columns}
\vspace{0.3cm}
\begin{alertblock}{Ý chính}
\texttt{print()} có thể in giá trị cố định và giá trị đang được lưu trong biến.
\end{alertblock}
\end{frame}

\begin{frame}[fragile]{In nhiều giá trị}
\texttt{print()} có thể nhận nhiều giá trị cùng lúc.
\begin{lstlisting}
name = "An"
age = 18
print("Name:", name)
print("Age:", age)
\end{lstlisting}
\begin{block}{Output}
\begin{verbatim}
Name: An
Age: 18
\end{verbatim}
\end{block}
\begin{itemize}
  \item Theo mặc định, Python chèn một dấu cách giữa các giá trị được in.
\end{itemize}
\end{frame}

\begin{frame}[fragile]{Ký tự phân cách mặc định}
\begin{lstlisting}
x = 10
y = 20
print(x, y)
\end{lstlisting}
\begin{block}{Output}
\begin{verbatim}
10 20
\end{verbatim}
\end{block}
\begin{center}
\begin{tikzpicture}[scale=0.95, every node/.style={transform shape}]
  \node[bluebox, minimum width=1.3cm] (a) at (0,0) {10};
  \node[box, minimum width=1.3cm] (space) at (2.0,0) {space};
  \node[greenbox, minimum width=1.3cm] (b) at (4.0,0) {20};
  \draw[arrow] (a) -- (space); \draw[arrow] (space) -- (b);
\end{tikzpicture}
\end{center}
\end{frame}

\begin{frame}[fragile]{Tham số \texttt{sep}}
\begin{block}{Mục đích}
\texttt{sep} điều khiển ký tự được đặt giữa các giá trị khi in.
\end{block}
\begin{columns}[T,onlytextwidth]
\column{0.48\textwidth}
\begin{lstlisting}
print("2026", "09", "06", sep="-")
\end{lstlisting}
\begin{verbatim}
2026-09-06
\end{verbatim}
\column{0.48\textwidth}
\begin{lstlisting}
print("Python", "Java", "C++", sep=" | ")
\end{lstlisting}
\begin{verbatim}
Python | Java | C++
\end{verbatim}
\end{columns}
\end{frame}

\begin{frame}[fragile]{Tham số \texttt{end}}
\begin{block}{Hành vi mặc định}
Sau mỗi lệnh \texttt{print()}, Python tự động xuống dòng.
\end{block}
\begin{columns}[T,onlytextwidth]
\column{0.48\textwidth}
\begin{lstlisting}
print("Hello")
print("Python")
\end{lstlisting}
\begin{verbatim}
Hello
Python
\end{verbatim}
\column{0.48\textwidth}
\begin{lstlisting}
print("Hello", end=" ")
print("Python")
\end{lstlisting}
\begin{verbatim}
Hello Python
\end{verbatim}
\end{columns}
\end{frame}

\begin{frame}[fragile]{Câu hỏi tự kiểm tra 1}
\begin{block}{Câu 1}
Output là gì?
\begin{lstlisting}
print("A", "B", "C", sep="-")
\end{lstlisting}
\end{block}
\begin{block}{Câu 2}
Output là gì?
\begin{lstlisting}
print("Hello", end=" ")
print("World")
\end{lstlisting}
\end{block}
\end{frame}

\begin{frame}{Đáp án tự kiểm tra 1}
\begin{enumerate}
  \item \texttt{A-B-C}
  \item \texttt{Hello World}
\end{enumerate}
\end{frame}

\section{Nhập dữ liệu với input()}

\begin{frame}[fragile]{Nhập dữ liệu với \texttt{input()}}
\begin{block}{Cú pháp}
\begin{lstlisting}
variable = input("Prompt")
\end{lstlisting}
\end{block}
\begin{exampleblock}{Ví dụ}
\begin{lstlisting}
name = input("Enter your name: ")
print("Hello", name)
\end{lstlisting}
\end{exampleblock}
\begin{alertblock}{Ý nghĩa}
Prompt được hiển thị cho người dùng. Chương trình chờ đến khi người dùng nhập dữ liệu và nhấn Enter.
\end{alertblock}
\end{frame}

\begin{frame}{\texttt{input()} làm gì?}
\begin{center}
\begin{tikzpicture}[scale=0.9, every node/.style={transform shape}]
  \node[bluebox, minimum width=2.7cm] (prompt) at (0,0) {Hiển thị\\lời nhắc};
  \node[orangebox, minimum width=2.7cm] (wait) at (3.8,0) {Chờ người dùng\\nhập dữ liệu};
  \node[purplebox, minimum width=2.7cm] (enter) at (7.6,0) {Người dùng\\nhấn Enter};
  \node[greenbox, minimum width=2.7cm] (return) at (11.4,0) {Trả về\\một chuỗi};
  \draw[arrow] (prompt) -- (wait);
  \draw[arrow] (wait) -- (enter);
  \draw[arrow] (enter) -- (return);
\end{tikzpicture}
\end{center}
\vspace{0.3cm}
\begin{block}{Bốn bước}
Hiển thị lời nhắc, chờ nhập, nhận dữ liệu vừa nhập và trả văn bản đó cho chương trình.
\end{block}
\end{frame}

\begin{frame}[fragile]{\texttt{input()} trả về chuỗi}
\begin{lstlisting}
age = input("Enter your age: ")
print(type(age))
\end{lstlisting}
Nếu người dùng nhập \texttt{18}, kết quả vẫn là:
\begin{verbatim}
<class 'str'>
\end{verbatim}
\begin{alertblock}{Quan trọng}
Ngay cả khi dữ liệu nhập trông giống một số, \texttt{input()} vẫn trả về dữ liệu kiểu \texttt{str}.
\end{alertblock}
\end{frame}

\section{Chuyển đổi kiểu dữ liệu đầu vào}

\begin{frame}{Vì sao cần chuyển kiểu?}
\begin{center}
\begin{tikzpicture}[scale=0.9, every node/.style={transform shape}]
  \node[bluebox, minimum width=3.0cm] (raw) at (0,0) {Dữ liệu từ\\bàn phím};
  \node[orangebox, minimum width=3.4cm] (str) at (4.0,0) {\texttt{input()} trả về\\\texttt{str}};
  \node[greenbox, minimum width=3.7cm] (num) at (8.4,0) {Chuyển kiểu trước\\khi tính toán};
  \draw[arrow] (raw) -- (str);
  \draw[arrow] (str) -- node[above]{\texttt{int()} hoặc \texttt{float()}} (num);
\end{tikzpicture}
\end{center}
\vspace{0.35cm}
\begin{block}{Quy tắc}
Khi dữ liệu nhập được dùng trong phép tính số học, hãy chuyển nó sang kiểu số phù hợp.
\end{block}
\end{frame}

\begin{frame}[fragile]{Chuyển sang số nguyên với \texttt{int()}}
\begin{lstlisting}
age = int(input("Enter your age: "))
print(age + 1)
\end{lstlisting}
Nếu người dùng nhập \texttt{18}, output là:
\begin{verbatim}
19
\end{verbatim}
\begin{itemize}
  \item Dùng \texttt{int()} cho số nguyên.
  \item Ví dụ: tuổi, số lượng, số sinh viên.
\end{itemize}
\end{frame}

\begin{frame}[fragile]{Chuyển sang số thực với \texttt{float()}}
\begin{lstlisting}
price = float(input("Enter price: "))
print(price)
\end{lstlisting}
\begin{itemize}
  \item Dùng \texttt{float()} cho số có thể chứa phần thập phân.
  \item Ví dụ: giá tiền, chiều cao, cân nặng, nhiệt độ.
\end{itemize}
\end{frame}

\begin{frame}[fragile]{Không chuyển kiểu}
\begin{lstlisting}
x = input("Enter x: ")
y = input("Enter y: ")
print(x + y)
\end{lstlisting}
Nếu người dùng nhập \texttt{10} và \texttt{20}, output là:
\begin{verbatim}
1020
\end{verbatim}
\begin{alertblock}{Vì sao?}
\texttt{x} và \texttt{y} là chuỗi, nên \texttt{+} thực hiện phép nối chuỗi.
\end{alertblock}
\end{frame}

\begin{frame}[fragile]{Có chuyển kiểu}
\begin{lstlisting}
x = int(input("Enter x: "))
y = int(input("Enter y: "))
print(x + y)
\end{lstlisting}
Output:
\begin{verbatim}
30
\end{verbatim}
\begin{alertblock}{Khác biệt}
Sau khi chuyển kiểu, \texttt{x} và \texttt{y} là số nguyên, nên \texttt{+} thực hiện phép cộng.
\end{alertblock}
\end{frame}

\begin{frame}[fragile]{Câu hỏi tự kiểm tra 2}
\begin{block}{Câu 1}
Đoạn code sau in gì nếu người dùng nhập \texttt{5}?
\begin{lstlisting}
x = input("x = ")
print(x * 3)
\end{lstlisting}
\end{block}
\begin{block}{Câu 2}
Cần sửa đoạn code như thế nào để kết quả là \texttt{15}?
\end{block}
\end{frame}

\begin{frame}[fragile]{Đáp án tự kiểm tra 2}
\begin{enumerate}
  \item Output là \texttt{555}. Giá trị là chuỗi \texttt{"5"}, nên nhân với 3 sẽ lặp lại chuỗi.
  \item Chuyển input sang số nguyên:
\end{enumerate}
\begin{lstlisting}
x = int(input("x = "))
print(x * 3)
\end{lstlisting}
\end{frame}

\section{Input - Process - Output}

\begin{frame}{Mô hình Input - Process - Output}
\begin{center}
\begin{tikzpicture}[scale=0.95, every node/.style={transform shape}]
  \node[bluebox, minimum width=3.4cm] (input) at (0,0) {Input\\\texttt{input()}};
  \node[orangebox, minimum width=3.4cm] (process) at (4.8,0) {Process\\tính toán};
  \node[greenbox, minimum width=3.4cm] (output) at (9.6,0) {Output\\\texttt{print()}};
  \draw[arrow] (input.east) -- (process.west);
  \draw[arrow] (process.east) -- (output.west);
\end{tikzpicture}
\end{center}
\begin{block}{Thói quen lập trình}
Với chương trình cho người mới bắt đầu, hãy xác định rõ input, xử lý và output trước khi viết code.
\end{block}
\end{frame}

\begin{frame}[fragile]{Ví dụ: diện tích hình chữ nhật}
\begin{lstlisting}
length = float(input("Length: "))
width = float(input("Width: "))
area = length * width
print("Area:", area)
\end{lstlisting}
\begin{columns}[T,onlytextwidth]
\column{0.32\textwidth}
\begin{block}{Input}
\texttt{length}, \texttt{width}
\end{block}
\column{0.32\textwidth}
\begin{block}{Process}
\texttt{length * width}
\end{block}
\column{0.32\textwidth}
\begin{block}{Output}
\texttt{print(...)}
\end{block}
\end{columns}
\end{frame}

\section{Định dạng output với f-string}

\begin{frame}[fragile]{Định dạng output với f-string}
f-string cung cấp cách rõ ràng để kết hợp text với giá trị của biến.
\begin{lstlisting}
name = "An"
age = 18
print(f"My name is {name}.")
print(f"I am {age} years old.")
\end{lstlisting}
\begin{block}{Output}
\begin{verbatim}
My name is An.
I am 18 years old.
\end{verbatim}
\end{block}
\end{frame}

\begin{frame}[fragile]{Biểu thức bên trong f-string}
\begin{lstlisting}
x = 10
y = 5
print(f"{x} + {y} = {x + y}")
\end{lstlisting}
Output:
\begin{verbatim}
10 + 5 = 15
\end{verbatim}
\begin{alertblock}{Ý chính}
Biểu thức trong \texttt{\{ \}} được tính trước khi chèn vào chuỗi.
\end{alertblock}
\end{frame}

\begin{frame}[fragile]{Định dạng số thực}
\begin{lstlisting}
price = 19.5678
print(f"Price: {price:.2f}")
\end{lstlisting}
Output:
\begin{verbatim}
Price: 19.57
\end{verbatim}
\begin{block}{Ý nghĩa}
\texttt{.2f} hiển thị số thực với 2 chữ số sau dấu thập phân.
\end{block}
\end{frame}

\section{Các lỗi thường gặp}

\begin{frame}[fragile]{Lỗi thường gặp: quên chuyển kiểu}
\begin{columns}[T,onlytextwidth]
\column{0.48\textwidth}
\begin{alertblock}{Sai}
\begin{lstlisting}
age = input("Age: ")
next_age = age + 1
\end{lstlisting}
\end{alertblock}
\column{0.48\textwidth}
\begin{exampleblock}{Đúng}
\begin{lstlisting}
age = int(input("Age: "))
next_age = age + 1
\end{lstlisting}
\end{exampleblock}
\end{columns}
\vspace{0.25cm}
\begin{block}{Nguyên nhân}
\texttt{age} là chuỗi, còn \texttt{1} là số nguyên. Không thể cộng trực tiếp hai kiểu này.
\end{block}
\end{frame}

\begin{frame}[fragile]{Lỗi thường gặp: input số không hợp lệ}
\begin{lstlisting}
age = int(input("Age: "))
\end{lstlisting}
\begin{block}{Vấn đề}
Nếu người dùng nhập \texttt{eighteen}, Python không thể chuyển chuỗi đó thành số nguyên.
\end{block}
\begin{alertblock}{Phạm vi bài học}
Hiện tại, giả sử người dùng nhập dữ liệu đúng định dạng yêu cầu. Kiểm tra dữ liệu nhập sẽ được học sau.
\end{alertblock}
\end{frame}

\begin{frame}[fragile]{Lỗi thường gặp: biến và chuỗi ký tự}
\begin{columns}[T,onlytextwidth]
\column{0.48\textwidth}
\begin{block}{In text \texttt{x}}
\begin{lstlisting}
x = 10
print("x")
\end{lstlisting}
Output:
\begin{verbatim}
x
\end{verbatim}
\end{block}
\column{0.48\textwidth}
\begin{block}{In giá trị của \texttt{x}}
\begin{lstlisting}
x = 10
print(x)
\end{lstlisting}
Output:
\begin{verbatim}
10
\end{verbatim}
\end{block}
\end{columns}
\end{frame}

\section{Ví dụ tổng hợp}

\begin{frame}[fragile]{Ví dụ tổng hợp 1: thông tin cá nhân}
\begin{lstlisting}
name = input("Name: ")
age = int(input("Age: "))

print(f"Hello {name}!")
print(f"Next year you will be {age + 1}.")
\end{lstlisting}
\begin{block}{Phân tích IPO}
\begin{itemize}
  \item Input: tên và tuổi
  \item Process: tính \texttt{age + 1}
  \item Output: lời chào và tuổi năm sau
\end{itemize}
\end{block}
\end{frame}

\begin{frame}[fragile]{Ví dụ tổng hợp 2: tính tổng tiền}
\begin{lstlisting}
price = float(input("Product price: "))
quantity = int(input("Quantity: "))
total = price * quantity

print(f"Total: {total:.2f}")
\end{lstlisting}
\begin{block}{Các chuyển đổi quan trọng}
\begin{itemize}
  \item \texttt{price}: dùng \texttt{float()}
  \item \texttt{quantity}: dùng \texttt{int()}
  \item \texttt{total}: hiển thị với \texttt{.2f}
\end{itemize}
\end{block}
\end{frame}

\begin{frame}[fragile]{Ví dụ tổng hợp 3: Celsius sang Fahrenheit}
Công thức:
\[
F = \frac{9}{5}C + 32
\]
\begin{lstlisting}
celsius = float(input("Celsius: "))
fahrenheit = 9 / 5 * celsius + 32
print(f"Fahrenheit: {fahrenheit:.2f}")
\end{lstlisting}
\end{frame}

\begin{frame}[fragile]{Dự đoán output: bài tập}
\begin{columns}[T,onlytextwidth]
\column{0.48\textwidth}
\begin{block}{Bài 1}
\begin{lstlisting}
x = "10"
y = "5"
print(x + y)
\end{lstlisting}
\end{block}
\begin{block}{Bài 2}
\begin{lstlisting}
x = 10
y = 5
print("Result:", x + y)
\end{lstlisting}
\end{block}
\column{0.48\textwidth}
\begin{block}{Bài 3}
\begin{lstlisting}
print("A", "B", sep=":", end=" ")
print("C")
\end{lstlisting}
\end{block}
\begin{block}{Bài 4}
\begin{lstlisting}
x = 7
print(f"x = {x}, x^2 = {x * x}")
\end{lstlisting}
\end{block}
\end{columns}
\end{frame}

\begin{frame}{Đáp án: dự đoán output}
\begin{enumerate}
  \item \texttt{105}
  \item \texttt{Result: 15}
  \item \texttt{A:B C}
  \item \texttt{x = 7, x\string^2 = 49}
\end{enumerate}
\end{frame}

\section{Bài tập thực hành}

\begin{frame}{Bài tập thực hành 1-4}
\begin{block}{Yêu cầu chung}
Viết trực tiếp các câu lệnh theo mô hình Input - Process - Output. Chưa dùng \texttt{def}.
\end{block}
\begin{enumerate}
  \item Lời chào: nhập tên và in \texttt{Hello, <name>!}
  \item Tuổi năm sau: nhập tuổi và in tuổi năm sau.
  \item Tổng và tích: nhập hai số nguyên, rồi in tổng và tích.
  \item Hình chữ nhật: nhập chiều dài và chiều rộng, rồi in diện tích và chu vi.
\end{enumerate}
\end{frame}

\begin{frame}{Bài tập thực hành 5-8}
\begin{enumerate}
  \item Tổng tiền mua hàng: nhập tên sản phẩm, đơn giá và số lượng, rồi in hóa đơn.
  \item Chuyển đổi nhiệt độ: đổi Celsius sang Fahrenheit và hiển thị 2 chữ số thập phân.
  \item Điểm trung bình: nhập ba điểm và hiển thị điểm trung bình với 2 chữ số thập phân.
  \item Chuyển đổi giây: nhập tổng số giây, rồi in số phút đầy đủ và số giây còn lại.
\end{enumerate}
\begin{block}{Toán tử hữu ích cho Bài 8}
\texttt{//} để chia lấy phần nguyên và \texttt{\%} để lấy phần dư.
\end{block}
\end{frame}

\begin{frame}{Bài thực hành tổng hợp: chi phí chuyến đi}
\begin{block}{Input}
Quãng đường đi được, mức tiêu hao nhiên liệu theo lít/100 km và giá nhiên liệu.
\end{block}
\begin{block}{Công thức}
\[
\text{fuel\_used} = \frac{\text{distance} \times \text{consumption}}{100}
\]
\[
\text{cost} = \text{fuel\_used} \times \text{fuel\_price}
\]
\end{block}
\end{frame}

\begin{frame}[fragile]{Bài thực hành tổng hợp: output mẫu}
\begin{verbatim}
Distance: 250
Consumption (L/100km): 7.5
Fuel price: 23000
Fuel used: 18.75 L
Estimated cost: 431250.00
\end{verbatim}
\begin{alertblock}{Trọng tâm}
Bài này kết hợp nhập dữ liệu, chuyển kiểu, tính toán và định dạng output.
\end{alertblock}
\end{frame}

\begin{frame}{Tự đánh giá sau bài học}
\small
\begin{itemize}
  \item Tôi biết sử dụng \texttt{print()}.
  \item Tôi hiểu \texttt{sep} và \texttt{end}.
  \item Tôi biết sử dụng \texttt{input()}.
  \item Tôi nhớ rằng \texttt{input()} trả về \texttt{str}.
  \item Tôi biết khi nào cần dùng \texttt{int()} hoặc \texttt{float()}.
  \item Tôi có thể viết chương trình theo mô hình Input - Process - Output.
  \item Tôi có thể định dạng số thực bằng \texttt{.2f}.
  \item Tôi có thể tự viết chương trình đơn giản mà chưa dùng \texttt{def}.
\end{itemize}
\end{frame}

\begin{frame}[fragile]{Tổng kết kiến thức}
Các công cụ quan trọng:
\begin{lstlisting}
print(...)
input(...)
int(...)
float(...)
type(...)
\end{lstlisting}
Các mẫu thường gặp:
\begin{lstlisting}
age = int(input("Age: "))
price = float(input("Price: "))
print(f"Result: {result:.2f}")
\end{lstlisting}
\end{frame}

\begin{frame}{Thông điệp cuối bài}
\begin{center}
\begin{tikzpicture}[scale=0.95, every node/.style={transform shape}]
  \node[bluebox, minimum width=3.2cm] (predict) at (0,0) {Predict};
  \node[greenbox, minimum width=3.2cm] (run) at (3.8,0) {Run};
  \node[orangebox, minimum width=3.2cm] (explain) at (7.6,0) {Explain};
  \node[purplebox, minimum width=3.2cm] (modify) at (3.8,-2.0) {Modify};
  \node[bluebox, minimum width=3.2cm] (practice) at (7.6,-2.0) {Practice};
  \draw[arrow] (predict) -- (run);
  \draw[arrow] (run) -- (explain);
  \draw[arrow] (explain.south) -- (practice.north);
  \draw[arrow] (practice) -- (modify);
  \draw[arrow] (modify.north) -- (run.south);
\end{tikzpicture}
\end{center}
\begin{alertblock}{Quy trình học}
Predict $\rightarrow$ Run $\rightarrow$ Explain $\rightarrow$ Modify $\rightarrow$ Practice.
\end{alertblock}
\end{frame}

\end{document}
